Let $\bitSequences := \lbrace 0, 1 \rbrace^{\N}$ be the infinite sequences of bits,
indexed from $0$,
and let $\neg$, $\wedge$, $\vee$ and $\oplus$ act on them position by position:
on a bit, $\neg f = 1 - f$, $\wedge$ is the minimum, $\vee$ the maximum,
and $\oplus$ the addition modulo $2$.
For $f \in \bitSequences$ and $n \in \N$,
write $\truncation{n}(f) := (f_0, \dots, f_n)$ for its first $n+1$ bits.

Two sets stand beside $\bitSequences$.
The first is $\truncatedSequences{n} := \Ftwo^{n+1}$, the tuples $(f_0, \dots, f_n)$,
read as polynomials of degree at most $n$ over the field $\Ftwo$ of two elements.
For $f \neq 0$ its \emph{degree} $\deg f$ is the largest $k$ with $f_k \neq 0$,
and $\deg 0 := -1$.
Position by position, $\oplus$ makes $\truncatedSequences{n}$ the abelian group
$(\Z/2\Z)^{n+1}$.
The second is $\remainders{n} := \Z/2^{n+1}\Z$, the remainders of the division by $2^{n+1}$,
with $\reduction{n} : \Z \to \remainders{n}$ the reduction taking the least non-negative
representative.
Notice that $\Ftwo = \remainders{0}$.

\begin{prop}\label{prop.binaryDigits}
 The map $\digitsOf{n} : \remainders{n} \to \truncatedSequences{n}$ sending
 $x \in \lbrace 0, \dots, 2^{n+1} - 1 \rbrace$ to its binary digits is a bijection.
\end{prop}
\begin{proof}
 The map $f \mapsto \sum_{k \leq n} f_k 2^k$ takes $\truncatedSequences{n}$ into
 $\lbrace 0, \dots, 2^{n+1} - 1 \rbrace$, and it is injective:
 at the highest $k$ where two tuples differ, their values differ by at least
 $2^k - \sum_{i < k} 2^i = 1$.
 Both sets have $2^{n+1}$ elements,
 so the map is a bijection, and $\digitsOf{n}$ is its inverse.
\end{proof}

The bijection is one of sets only.
Transported through $\digitsOf{n}$, the addition of $\remainders{n}$ makes
$\truncatedSequences{n}$ the cyclic group of order $2^{n+1}$,
which is not the group of $\oplus$ as soon as $n \geq 1$:
$1 + 1 = 2$ carries a digit, where $1 \oplus 1 = 0$ does not.

Identifying a tuple of $\truncatedSequences{n}$ with its extension by zeros,
write $\eventuallyZero := \bigcup_{n} \truncatedSequences{n} \subset \bitSequences$ for the
\emph{eventually zero} sequences, those with $f_k = 0$ for every $k$ beyond some $n$,
the least such $n$ being $\deg f$;
and let $\binaryValue : \eventuallyZero \to \N$, $\binaryValue(f) := \sum_{k} f_k 2^k$,
the sum being finite.

\begin{prop}\label{prop.binaryRepresentation}
 $\binaryValue$ is a bijection.
 We call its inverse $\binaryRepresentation : \N \to \eventuallyZero$ the binary representation.
\end{prop}
\begin{proof}
 Two eventually zero sequences lie in a common $\truncatedSequences{n}$,
 on which $\binaryValue$ is injective by Proposition \ref{prop.binaryDigits};
 and every $m \in \N$ lies below $2^{n+1}$ for some $n$,
 so it is the value of a sequence of $\truncatedSequences{n}$, by the same proposition.
\end{proof}

The two are related by truncation.

\begin{prop}\label{prop.truncationCommutes}
 $\truncation{n} \circ \binaryRepresentation = \digitsOf{n} \circ \reduction{n}$ on $\N$.
\end{prop}
\begin{proof}
 Write $\binaryRepresentation(m) = (f_k)_k$ and split the sum at $n$:
 $m = \sum_{k \leq n} f_k 2^k + \sum_{k > n} f_k 2^k$.
 The second sum is a multiple of $2^{n+1}$ and the first lies in
 $\lbrace 0, \dots, 2^{n+1} - 1 \rbrace$, so the first is $\reduction{n}(m)$,
 and its digits are $\truncation{n}(\binaryRepresentation(m))$.
\end{proof}
